\section{Data preparation}
\label{sec:preparacion-datos}

\subsection{Consolidation and cleaning}

The six annual SIMA files do not share a structure: the headers, the
position of the units row and the number of stations change (13 in 2020,
15 from 2022 onward). They are consolidated into a single hourly table of
15 parameters by station and date, 754,603 records between 2020 and 2025.
The quality audit detects sentinel values ($-9999$, the logger's missing-data
code) and physically impossible values (relative humidity of up to
725\,\%, temperatures of 785\,°C), which are set to missing rather than
imputed; details are given in the data dictionary (\cref{sec:dic:centinela}).
The share of missing data per parameter ranges from 3.66\,\% for solar
radiation to 25.23\,\% for \ce{PM_{2.5}}, but it is not evenly distributed
across years or stations, which determines the study window.

\subsection{Data selection}

The dataset for the subsequent steps is built from the data dictionary and
the quality diagnosis.

Since the objective is to analyze the behavior of tropospheric ozone, each
variable is assessed to determine, based on the literature, whether there is
a documented mechanism ---chemical, physical or meteorological--- through
which it influences \ce{O3} formation.

The target variable is the \ce{O3} from the database:

\begin{table*}[pos=t]
  \centering
  \caption{Target variable: tropospheric ozone (\ce{O3}).}
  \label{tab:variable-objetivo}
  \begin{tabular}{lcccc}
    \toprule
    Unit & Operating range & Observed range & Median & Missing (\%) \\
    \midrule
    ppb & 0--153 to 0--185 & 0.5--265.0 & 23.00 & 14.73 \\
    \bottomrule
  \end{tabular}
\end{table*}

\begin{quote}
    Stations report results every hour, i.e., 24 measurements per day. The
    regulatory MDA8 metric collapses those 24 measurements into a single
    value per station and day, since NOM-020-SSA1 sets the ozone limit on
    the 8-hour average, not on the hourly value.
\end{quote}

\Cref{tab:nom-limites} lists the current Mexican limit values for ozone and
particulate matter. Note that the regulatory indicator for \ce{O3} is not the
hourly concentration but the 8-hour moving average, and that the evaluation
criterion is the maximum of the daily maxima of those averages; this
justifies adopting the daily maximum 8-hour moving average (MDA8) as the
target variable of the analysis.

\begin{table*}[pos=t]
  \centering
  \caption{Current Mexican air quality limit values for the pollutants of
    interest. The years indicate the gradual compliance schedule set by each
    standard \parencite{nom020ssa1, nom025ssa1}.}
  \label{tab:nom-limites}
  \footnotesize
  \begin{tabular}{p{2.6cm} p{7.2cm} ccc}
    \toprule
    & & \multicolumn{3}{c}{Limit value} \\
    \cmidrule(lr){3-5}
    Indicator & Evaluation criterion & Year 1 & Year 3 & Year 5 \\
    \midrule
    \multicolumn{5}{l}{\textbf{\ce{PM10}} ($\mu$g/m$^{3}$) --- NOM-025-SSA1-2021, DOF 27 Oct 2021} \\
    \addlinespace[2pt]
    24 hours & 99th percentile (tolerated frequency: 1\,\% of the time per year) & 70 & 60 & 50 \\
    Annual   & Annual mean of the 24-hour averages                              & 36 & 28 & 20 \\
    \midrule
    \multicolumn{5}{l}{\textbf{\ce{PM_{2.5}}} ($\mu$g/m$^{3}$) --- NOM-025-SSA1-2021, DOF 27 Oct 2021} \\
    \addlinespace[2pt]
    24 hours & 99th percentile (tolerated frequency: 1\,\% of the time per year) & 41 & 33 & 25 \\
    Annual   & Annual mean of the 24-hour averages                              & 10 & 10 & 10 \\
    \midrule
    \multicolumn{5}{l}{\textbf{Ozone} (\ce{O3}, ppm) --- NOM-020-SSA1-2021, DOF 28 Oct 2021} \\
    \addlinespace[2pt]
    1 hour                 & Maximum of the daily maxima of the hourly concentrations                  & 0.090 & 0.090 & 0.090 \\
    8-hour moving average  & Maximum of the daily maxima of the 8-hour moving average concentrations   & 0.065 & 0.060 & 0.051 \\
    \bottomrule
  \end{tabular}
\end{table*}

Missing \ce{O3} values by station and year (\cref{sec:dic:faltantes}) show
that 2020 is unusable: of the 13 stations operating that year, 7 have between
89\,\% and 100\,\% missing data. The year 2020 is discarded entirely.

In 2021 missing data remain high at several stations, but within a range
that allows the analysis; from 2022 onward no station exceeds 20\,\%, except
\texttt{NO3} in 2025 (30.3\,\%). The full years 2021 to 2025 are therefore
retained: 1,826 calendar days.

Within this window, no station has enough missing data to justify excluding
it, so the 15 levels of the variable \texttt{estacion} (station) are kept.
This variable is not subject to the same criterion as the measured
parameters: it is not a physicochemical quantity but an identifier that
defines the spatial structure of the dataset.

Next, the variables that affect tropospheric ozone levels are identified.
Information on the chemical, physical and meteorological processes that
promote or hinder \ce{O3} formation is compiled from the literature.

\begin{itemize}
    \item \textbf{Solar radiation}: in summer, sunlight is more intense and
      temperatures rise, which accelerates photochemical reactions and
      produces higher ozone levels.
    \item \textbf{Quantified seasonal variation}: mean concentrations of
      $45.6\,\mu\mathrm{g}/\mathrm{m}^{3}$ in summer,
      $47.3\,\mu\mathrm{g}/\mathrm{m}^{3}$ in spring,
      $38.0\,\mu\mathrm{g}/\mathrm{m}^{3}$ in autumn and
      $33.6\,\mu\mathrm{g}/\mathrm{m}^{3}$ in winter.
    \item \textbf{Diurnal variation}: ozone levels tend to peak in the late
      afternoon and decrease at night due to the lack of sunlight and to
      titration by \ce{NO}, which consumes ozone.
    \item \textbf{Topography}: valleys surrounded by mountains trap pollutants
      and favor ozone formation (stagnation effect), which applies in
      particular to the Monterrey metropolitan area.
    \item \textbf{Acute events}: heat waves cause atypical short-term
      increases.
\end{itemize}

The chemistry section of Sillman's article \parencite{SanfordSillman}
describes the formation mechanism of tropospheric ozone:

\begin{quote}
    Ozone formation in the lower atmosphere results from a series of
    photochemical reactions involving precursors emitted by natural and
    anthropogenic sources. These reactions begin with the photolysis of
    nitrogen dioxide (\ce{NO2}) under sunlight, which produces atomic oxygen
    (\ce{O}). This atomic oxygen then reacts with molecular oxygen (\ce{O2})
    to form ozone (\ce{O3}).
\end{quote}

The complete cycle:

\begin{enumerate}
    \item $\ce{NO2} + h\nu \to \ce{NO}+\ce{O}$ (photolysis)
    \item $\ce{O} + \ce{O2} + \ce{M} \to \ce{O3} + \ce{M}$ (ozone formation)
    \item VOCs (volatile organic compounds) react with free radicals to
      generate additional \ce{NO2}, feeding the cycle continuously.
\end{enumerate}

From this information, the meteorological variables used in the rest of the
analysis are determined.

\begin{itemize}
    \item \texttt{TOUT} (ambient temperature): the higher the temperature,
      the faster ozone forms.
    \item \texttt{RAINF} (precipitation): precipitation implies cloudy skies,
      which reduce exposure to solar radiation.
    \item \texttt{PRS} (atmospheric pressure): high pressure brings clear,
      cloudless days and near-calm winds, and also acts as a \textit{lid}:
      the heavy air sinks, compresses pollutants against the ground and
      prevents their dispersion.
    \item \texttt{WSR} and \texttt{WDR} (wind speed and direction):
      determine the ability of the atmosphere to disperse pollutants.
\end{itemize}

Among the air pollutants:

\begin{itemize}
    \item \ce{CO} (carbon monoxide) acts as a proxy for the VOCs that react
      with free radicals.
    \item \ce{NO2} is the main ingredient of the photochemical reaction with
      solar radiation.
    \item \ce{NOx} is the group $\ce{NO2} + \ce{NO}$, essential to \ce{O3}
      formation.
    \item \texttt{PM10} and \texttt{PM2.5} are not gases but solid particles
      suspended in the air; high concentrations act as a screen against
      sunlight and attenuate the reactions.
\end{itemize}

\subsection{Daily aggregation}
\label{sec:agregacion-diaria}

The unit of analysis of the models is the station-day, because the standard
is evaluated on the MDA8 and not on the hour. The clean hourly table is
collapsed to one row per day and station with 16 continuous summaries: the
precursors \ce{NO}, \ce{NO2}, \ce{NOx} and \ce{CO} as the mean of the
06--09 h window; maximum temperature, cumulative radiation, mean and
daytime-minimum humidity, mean pressure, mean and minimum wind speed, mean
vector wind components, and daily means of \ce{PM10}, \ce{PM_{2.5}} and
\ce{SO2}. The MDA8 is declared valid only if at least 18 of the 24 8-hour
moving windows of the day are complete; of the 26,691 station-day rows,
25,100 meet this rule. Each summary and the completeness rule are justified
in the aggregation log of the repository. Each station-day is assigned its
\texttt{zona} (zone), \texttt{temporada} (season) and \texttt{tipo\_dia}
(day type), which distinguishes weekdays, Saturdays, Sundays and the 81
public holidays of the period, whose traffic pattern is anomalous.
